% Generated draft: requires mathematical and source review.
\documentclass[11pt]{article}
\usepackage[margin=25mm]{geometry}
\usepackage{fontspec}
\setmainfont{Latin Modern Roman}
\usepackage{amsmath,amssymb,mathtools,unicode-math}
\setmathfont{Latin Modern Math}
\usepackage{enumitem,needspace,xurl,hyperref}
\setlength{\parindent}{0pt}
\setlength{\parskip}{5pt}
\setlength{\emergencystretch}{3em}
\newcommand{\N}{\mathbb{N}}
\newcommand{\Z}{\mathbb{Z}}
\newcommand{\Q}{\mathbb{Q}}
\newcommand{\R}{\mathbb{R}}
\newcommand{\C}{\mathbb{C}}
\newcommand{\F}{\mathbb{F}}
\newcommand{\A}{\mathbb{A}}
\newcommand{\PP}{\mathbb{P}}
\newcommand{\E}{\mathbb{E}}
\newcommand{\eps}{\varepsilon}
\newcommand{\dd}{\,\mathrm{d}}
\begin{document}
\section*{UnsolvedMath ID: 800638}
\textbf{Source problem code:} catalogue entry
\subsection{Definitions and mathematical statement}
The supplied question is
\[
\mathsf{add}\bigl(\mathsf{U}_{\mathrm{fin}}(\Gamma,\mathcal O)\bigr)=\mathfrak b\;?
\]
The record does not specify the underlying class of spaces or its cover conventions. The following gives a conditional formulation using countable open covers; these conventions require confirmation against the source.

For a topological space \(Y\), let \(\mathcal O(Y)\) denote the countable families of open subsets of \(Y\) that cover \(Y\), with \(Y\) itself excluded as a member. Let \(\Gamma(Y)\) consist of the countably infinite members \(\mathcal U\) of \(\mathcal O(Y)\) such that each point of \(Y\) belongs to all but finitely many members of \(\mathcal U\). The property \(\mathsf{U}_{\mathrm{fin}}(\Gamma,\mathcal O)\) means that for every sequence \((\mathcal U_n)_{n\in\N}\) from \(\Gamma(Y)\), each having no finite subcover, there exist finite subfamilies \(\mathcal F_n\subseteq\mathcal U_n\) such that
\[
\left\{\bigcup\mathcal F_n:n\in\N\right\}\in\mathcal O(Y).
\]
Write \(\mathcal P\) for this property.

To specify the union operation in the additivity number, fix an ambient topological space \(E\) and give every subset of \(E\) its subspace topology. Define, whenever a witnessing family exists,
\[
\mathsf{add}_E(\mathcal P)
=\min\left\{|I|:\ X_i\subseteq E\text{ and }X_i\in\mathcal P\ (i\in I),\quad \bigcup_{i\in I}X_i\notin\mathcal P\right\}.
\]
Here \(X\in\mathcal P\) means that \(X\) has the stated selection property. For \(f,g\in\N^{\N}\), write \(f\leq^*g\) if \(f(n)\leq g(n)\) for all but finitely many \(n\). The bounding number is
\[
\mathfrak b=\min\left\{|B|:B\subseteq\N^{\N},\quad \forall g\in\N^{\N}\ \exists f\in B\ (f\not\leq^*g)\right\}.
\]
The intended problem asks whether the additivity number of \(\mathcal P\) equals \(\mathfrak b\). Under the conditional ambient-space formulation above, its target is \(\mathsf{add}_E(\mathcal P)=\mathfrak b\) for the intended ambient space and admissible subspaces. The supplied record does not identify these, so this is not an assertion or question quantified over all ambient spaces \(E\). Any restrictions on admissible subspaces must also be incorporated into the definition of additivity.
\subsection{Short English statement}
Does the least number of spaces with the selection property \(\mathsf{U}_{\mathrm{fin}}(\Gamma,\mathcal O)\) whose union fails that property equal the bounding number \(\mathfrak b\)? The intended class of spaces and cover conventions remain unspecified in the supplied record.
\subsection{Sources}
\begin{itemize}
\item Local UnsolvedMath registry, permanent ID \(800638\), exact source problem code \texttt{catalogue entry}, title ``Selection Principles Milestone Problems --- Problem 10 ([AddQuad]).'' The record attributes the problem to Tsaban, with proposed year 2003.
\item Supplied source: Tsaban, \emph{Selection Principles in Mathematics: A Milestone of Open Problems} (2003), Problem 10. Supplied URL: \url{https://arxiv.org/abs/math/0312182}. The record reports access on 2026-10-16 and extraction by \texttt{arxiv\_latex}; these are supplied provenance details, not an independently verified visit.
\item The source excerpt places this question within the general problem of expressing additivity numbers of selection properties through cardinal characteristics such as \(\mathfrak b\) and \(\mathfrak d\), and describes it as one of three particularly interesting cases. It refers to [LecceSurvey] for a list of properties with open additivity questions. No bibliographic identification for [LecceSurvey] is supplied.
\item The supplied attribution/context label is [AddQuad]; its bibliographic identification is not supplied. The registry labels this problem ``open''; that label is supplied curation, not an independently verified assessment of its present status.
\end{itemize}
\end{document}
